% CV — Tran Quoc Vy
% Rebuilt 2026-11-17 to match the original layout (Computer Modern, single page).
% Build: tectonic CV_TranQuocVy_BackendEngineer.tex   (or upload to Overleaf, pdfLaTeX)
% Every project claim is backed by merged PRs; see CONTRIBUTIONS_AUDIT.md (kept outside this repo) before editing.

\documentclass[11pt,a4paper]{article}

\usepackage[T1]{fontenc}
\usepackage[utf8]{inputenc}
\usepackage[a4paper,top=1.1cm,bottom=1.0cm,left=1.5cm,right=1.5cm]{geometry}
\usepackage{titlesec}
\usepackage{enumitem}
\usepackage[hidelinks]{hyperref}
\usepackage{tabularx}
\usepackage{xcolor}

% Phone number lives in phone.local.tex (gitignored): \def\CVPhone{...}
\IfFileExists{phone.local.tex}{\input{phone.local.tex}}{}

\pagestyle{empty}
\setlength{\parindent}{0pt}
\setlength{\parskip}{0pt}

% Section headings: bold uppercase with a full-width rule underneath
\titleformat{\section}{\large\bfseries}{}{0em}{\MakeUppercase}[\vspace{-7pt}\rule{\linewidth}{0.6pt}]
\titlespacing*{\section}{0pt}{6pt}{3pt}

\setlist[itemize]{leftmargin=1.1em, label=\textbullet, itemsep=0.5pt, topsep=1pt, parsep=0pt}

% \entry{Title}{Date}{Subtitle}
\newcommand{\entry}[3]{%
  \begin{tabularx}{\linewidth}{@{}X r@{}}
    \textbf{#1} & #2 \\
  \end{tabularx}\\[-2pt]
  {\itshape #3}\par
}

\begin{document}

% ---------- Header ----------
\begin{center}
  {\fontsize{24}{28}\selectfont TRAN QUOC VY}\\[3pt]
  {\large\bfseries BACKEND ENGINEER INTERN}\\[4pt]
  {\small \href{mailto:tqvy.work@gmail.com}{tqvy.work@gmail.com} \textbar{}
  \ifdefined\PublicVersion\else\ifdefined\CVPhone \CVPhone{} \textbar{} \fi\fi
  \href{https://tranquocvy.github.io}{tranquocvy.github.io} \textbar{}
  \href{https://github.com/tranquocvy}{github.com/tranquocvy} \textbar{}
  \href{https://www.linkedin.com/in/tqvy/}{linkedin.com/in/tqvy}}
\end{center}
\vspace{-4pt}

% ---------- Summary ----------
\section{Summary}
Final-year Software Engineering student at HCMUS focused on backend engineering. In team projects I own the
parts that must stay correct when things fail: payments with idempotency and circuit breakers, message-queue
workers, Redis caching and CI/CD. Case studies with links to my merged PRs:
\href{https://tranquocvy.github.io}{tranquocvy.github.io}.

% ---------- Education ----------
\section{Education}
\entry{University of Science (HCMUS) -- VNU-HCM}{2023 -- 2027}{Bachelor of Software Engineering}
\begin{itemize}
  \item \textbf{GPA:} 8.83/10 (3.788/4.0)
  \item \textbf{Relevant Coursework:} Data Structures and Algorithms, Object-Oriented Programming, Database Systems,
        Web Application Development, Software Engineering, Software Design, Introduction to DevOps
\end{itemize}

% ---------- Skills ----------
\section{Skills}
\begin{itemize}
  \item \textbf{Backend:} Node.js, NestJS, Express, REST API design, JWT authentication, role-based authorization
  \item \textbf{Data \& Messaging:} PostgreSQL, Prisma, Knex, Redis, RabbitMQ, MongoDB (basic)
  \item \textbf{DevOps:} Docker, GitHub Actions, Git/Gitflow, Linux, Railway; Jenkins and Kubernetes (basic)
  \item \textbf{Practices:} Unit testing, design patterns (Strategy, cache-aside, circuit breaker), code review via pull requests, Agile/Scrum
  \item \textbf{Languages:} TypeScript, JavaScript, C++, C\#
\end{itemize}

% ---------- Projects ----------
\section{Projects}
\entry{\href{https://tranquocvy.github.io/projects/ticket-box/}{Ticket\_Box -- Event Ticketing Platform}}{May 2026 -- Jul 2026}%
{Backend lead for payments \& background jobs, team of 4 \textbullet{} 17 merged PRs, 41\% of commits}
\begin{itemize}
  \item \textbf{Payment reliability:} Built the PayOS payment module behind a Strategy interface with two-layer
        idempotency (Redis \texttt{SET NX} reservation plus a unique key in PostgreSQL) to stop duplicate charges.
  \item \textbf{Fault tolerance:} Wrapped gateway calls in an opossum circuit breaker (3\,s timeout, opens at 50\%
        errors) that fails fast with \texttt{503}; reconciled timed-out transactions and released held seats on
        cancellation or failed webhooks.
  \item \textbf{Async workers:} Built RabbitMQ consumers for streamed CSV guest-list import (batched PostgreSQL
        upserts with progress tracking), AI artist bios from PDF press kits via Gemini, and email/in-app notifications.
  \item \textbf{Performance \& delivery:} Added Redis cache-aside for concert reads and a token-bucket rate limiter;
        set up GitHub Actions CI (lint, typecheck, Prisma validation, migration check) and Docker deployment to Railway.
  \item \textbf{Stack:} NestJS, Prisma, PostgreSQL, Redis, RabbitMQ, PayOS, Docker; unit tests for payment timeout,
        duplicate-payment and circuit-breaker paths.
\end{itemize}

\vspace{3pt}
\entry{\href{https://tranquocvy.github.io/projects/retrobit/}{Retrobit -- Retro Board Game Platform}}{Jan 2026}%
{Backend owner, team of 4 \textbullet{} Express, PostgreSQL, Knex, JWT, React, Vite}
\begin{itemize}
  \item \textbf{API foundation:} Designed and built the Express/PostgreSQL REST API for users, games and
        resumable game sessions, with Knex migrations and seed data.
  \item \textbf{Auth \& security:} Implemented JWT authentication, role-based authorization middleware and email
        OTP verification over SMTP.
  \item \textbf{Rating system:} Wrote the ELO engine: dynamic K-factor, fixed bot ratings for player-vs-AI modes,
        and score-to-result mapping for single-player games; plus the achievements service.
  \item \textbf{Delivery:} Shipped the Match-3 and Memory games on the React frontend and deployed both tiers;
        opened 23 PRs, the most on the team.
\end{itemize}

\end{document}
